% -*- root: ../thesis.tex -*-
%!TEX root = ../thesis.tex
% ******************************************************************************
% ************************* Custom LaTeX Commands ******************************
% ******************************************************************************
%
% This file contains custom LaTeX commands that may be used across chapters.
% Add your command definitions below.

% Example:
% \newcommand{\mycommand}[1]{\textbf{#1}}

%%% Chapter engram dynamics
% Reusable math commands migrated from the manuscript preamble.
\DeclareMathOperator{\diag}{diag}
\newcommand{\var}[1]{\text{Var}\left[#1\right]}
\newcommand{\mc}{\mathcal}
\newcommand{\bd}{\boldsymbol}
\newcommand{\mb}{\mathbf}
\newcommand{\dg}{\dagger}
\newcommand{\ddt}{\frac{\mathrm{d}}{\mathrm{d}t}}

%%% Chapter PPT
% commands for outside math mode
\newcommand{\ppca}{\texorpdfstring{PP$_\textrm{\footnotesize CA1}$}{PP-CA1}}
\newcommand{\ppsub}{\texorpdfstring{PP$_\textrm{\footnotesize SUB}$}{PP-SUB}}
% commands for inside math mode (in equationa, eq)
\newcommand{\ppcaeq}{\rm PP\text{-}CA1}
\newcommand{\ppsubeq}{\rm PP\text{-}SUB}
\newcommand{\casubeq}{\rm CA1\text{-}SUB}

\renewcommand{\vec}[1]{\mathbf{#1}} % Michiel added this to change to bold fonts for vectors (looks cleaner)
\newcommand{\vc}[1]{\vec{#1}}

\newcommand{\W}{\vc{W}} % weights of PP
\newcommand{\V}{\vc{V}} % weights of SC
\newcommand{\x}{\vc{x}} % input rates from EC
\newcommand{\z}{\vc{x}'} % input rates from CA3
\newcommand{\y}{y} % output rate of the CA1 cell
\newcommand{\LW}{L} % learning window
\newcommand{\e}{{\rm e}} %
\newcommand{\Cxx}{\langle \x \x^{\mathsf{T}}\rangle} % covariance matrix of the direct input
\newcommand{\Cxz}{\langle \x \z^{\mathsf{T}}\rangle} % covariance between direct and indirect input
\newcommand{\HS}[1]{{\color{blue} #1}}

%%% FIGURE MACROS %%%
% The variables here are names after the original figure 1-4

%% Figure 1 %%
% Figure 1 title
\newcommand{\figureOneFigureTitle}{%
  A mechanistic basis for systems memory consolidation%
}

% Figure 1 label
\newcommand{\figureOneFigureLabel}{fig:PPT_mechanism}

% Figure 1 caption
\newcommand{\figureOneFigureCaption}{%
  (A) Circuit motif for the parallel pathway theory.
  Cue-response associations are initially stored in an indirect synaptic pathway (blue)
  and consolidated into a parallel direct pathway (red).
  (B) Hippocampal connectivity.
  The entorhinal cortex projects to CA1 through an indirect pathway via DG-CA3 and the Schaffer collaterals (SC, blue arrow), and through the direct perforant path (\ppca, red arrow).
  (C) Model of consolidation through STDP.
  Left: before consolidation, a strong SC input (middle, blue vertical bar) causes a large EPSP and triggers a spike in CA1 (bottom, black vertical bar).
  A weak \ppca~input (top, red) that precedes the SC input is potentiated by STDP.
  Right: after consolidation through STDP, the \ppca~input (top) can trigger a spike in CA1 by itself (bottom).
  (D-E) Consolidation in a single integrate-and-fire CA1 cell receiving 1000 \ppca~and 1000 SC excitatory inputs.
  (D) \ppca~activity consists of independent poisson spike trains; the SC activity is an exact copy of the \ppca~activity, delayed by 5 ms.
  (E) Consolidation of a synaptic weight pattern from non-plastic SC synapses to plastic \ppca~synapses.
  Left and middle: normalized synaptic weights before and after consolidation.
  Right: time course of correlation between SC and \ppca~weight vectors during consolidation (mean $\pm$ SEM for 10 trials).
  \added{(F) 
  Failure of consolidation of a synaptic weight pattern from non-plastic PP$_\mathrm{CA1}$ to plastic SC synapses; panels as in E. 
}}

% Figure 1 filename
%\newcommand{\figureOneFigureFilename}{Fig_1_with_motif.pdf}
\newcommand{\figureOneFigureFilename}{fig1}


% Figure 2 title
\newcommand{\figureTwoFigureTitle}{%
  Consolidation of spatial representations%
}

% Figure 2 label
\newcommand{\figureTwoFigureLabel}{fig:spatial}

% Figure 2 caption
\newcommand{\figureTwoFigureCaption}{%
  (A) Replay of \ppca~and SC activity during sleep.
  500 \ppca~inputs and 2500 SC inputs are spatially tuned on a linear track with periodic grid fields (top, red) and place fields (bottom, blue).
  Spiking activities are independent Poisson processes (10 spikes/s) inside place/grid fields, otherwise silent.
  SC activity is delayed by 5 ms.
  (B) Multi-compartmental model of a reconstructed CA1 pyramidal neuron (see %\hyperref[ssec:methods_multicompartment]{
  {\em Methods}).
  \ppca~and SC inputs project to distal apical tuft dendrites (red dots) and proximal apical and basal dendrites (blue dots).
  (C) Active neuron properties.
  Top: somatic sodium spike (black) propagates to the distal tuft and initiates a dendritic calcium spike (red) and further sodium spikes.
  Bottom: dendritic calcium spike leads to bursts of somatic spikes.
  (D) Spatial tuning before consolidation.
  SC provides place field-tuned input to the CA1 cell (left, blue), which yields spatially tuned spiking activity (right, blue); \ppca~input is not spatially tuned (left, red), and (alone) triggers low and untuned spiking activity (right, red).
  (E) Somatic and dendritic activity during consolidation.
  During replay, SC input generates backpropagating sodium spikes (black vertical lines) that generate dendritic calcium spikes (red).
  (F) After consolidation.
  Spatial tuning is consolidated from the indirect SC pathway into the direct \ppca~pathway.
  Left: spatial tuning of total \ppca~input (red) approaches theoretically derived \ppca~input tuning (magenta; see {\em Methods}).
  Right: CA1 output is place field-tuned through either SC or \ppca~input alone.
  (G) Evolution of correlation between actual and optimal \ppca~input tuning (see F) for replay speeds corresponding to hippocampal replay events (black) and real-time physical motion (grey). Position in D, E, and F normalized to [0,1].
}  % end figureTwoFigureCaption

% Figure 2 filename
%\newcommand{\figureTwoFigureFilename}{Fig_2.pdf}
\newcommand{\figureTwoFigureFilename}{fig3}


% Figure 3 title
\newcommand{\figureThreeFigureTitle}{%
  Consolidation of place-object associations in multiple hippocampal stages%
}

% Figure 3 label
\newcommand{\figureThreeFigureLabel}{fig:lesion}

% Figure 3 caption
\newcommand{\figureThreeFigureCaption}{%
  (A) Structure of the extended model.
  \ppsub: perforant path to the subiculum.
  Each area (EC, DG-CA3, CA1, SUB) contains object-coding and place-coding populations.
  Open arrows: all-to-all connections between these areas.
  (B) Decoding of consolidated associations.
  Top: The location of a platform in a circular environment is stored as an object-place association in the SC (thick diagonal arrows in A, right).
  Middle: Platform position probability maps given the platform object cue, inferred from the CA1 output resulting from SC or \ppca~alone, at different times during consolidation ({\em Methods}).
  Bottom: Platform-in-quadrant probabilities ($
 \pm \mathrm{SEM}$) given \ppca~input alone during consolidation.
  Quadrant with correct platform position (target quadrant) in orange.
  (C) Consolidation from SC to \ppca~and to \ppsub~over four weeks.
  Each day, a new association is first stored in SC and then partially consolidated.
  An association on day~0 is monitored in SC, \ppca,~and \ppsub.
  Panels as in B.
  (D) Effects of \ppca~lesions on memory consolidation, model and experiment (data with permission from \cite{Remondes2004}).
  Histograms of time ($
 \pm \mathrm{SEM}$) spent in quadrants at different delays after memory acquisition (``probe'').
  Dashed lines at 25\% are chance levels.
  T: target quadrant; Left, Right: adjacent quadrants; O: opposite quadrant.
  Top: Control without lesion.
  Middle: Lesion before memory acquisition.
  Bottom: Lesion 21 days after memory acquisition.
}

% Figure 3 filename
%\newcommand{\figureThreeFigureFilename}{Fig_3.pdf}
\newcommand{\figureThreeFigureFilename}{fig4}


% Figure 4 title
\newcommand{\figureFourFigureTitle}{%
  Consolidation from hippocampus into neocortex by a hierarchical nesting of consolidation circuits%
}

% Figure 4 label
\newcommand{\figureFourFigureLabel}{fig:hierarchical}

% Figure 4 caption
\newcommand{\figureFourFigureCaption}{%
  (A) Schematic of the hierarchical model.
  The hippocampal formation (HPC) is connected to cortical input circuit 1 and output circuit 1.
  Increasing numbers indicate circuits further from the HPC and closer to the sensory/motor periphery.
  Each direct connection at one level (e.g., dark blue arrow between input 1 and output 1) is part of the indirect pathway of the next level (e.g., for pathways from input 2 to output 2).
  Learning rates of the direct connections decrease exponentially with increasing level (i.e., from blue to red).
  (B) Memories gradually propagate to circuits more distant from the HPC.
  The correlation of the initial HPC weights with the direct pathways is shown as a function of time and reveals a memory wave from HPC into neocortex.
  The maximum of the output circuits follows approximately a power-law (black curve).
  Noise level indicates chance-level correlations between pathways.
  (C) Consolidated memories yield faster responses (from sensory periphery, e.g., Input 8, to system output) because these memories are stored in increasingly shorter synaptic pathways.
}

% Figure 4 filename
%\newcommand{\figureFourFigureFilename}{Fig_4.pdf}
\newcommand{\figureFourFigureFilename}{fig5}


% This is copied from the PLOS Latex template: Includes just the caption, not the figure
\newcommand{\figureCaptionOnly}[5][!ht]{%
    % Place figure captions after the first paragraph in which they are cited.
    \begin{figure}[#1]
    \caption{{\bf #2.} #4}
    \label{#3}
    \end{figure}
}

\newcommand{\makeWideFigure}[5][!htbp]{%
  % turn off line numbers in captions

  \begin{figure}[#1]
  \begin{adjustwidth}{-2.25in}{0in}
  \includegraphics[width=\linewidth]{#5}

  % wide figure with multi-page caption
  \vspace{.4cm} % set vertical space between text and figure
  \end{adjustwidth}
  \caption{\textbf{#2.} #4}
  \label{#3}
  \end{figure}
}

\newcommand{\makeStandardFigure}[5][H]{%
  % standard floating figure

  \begin{figure}[#1]
    % use to correct figure counter if necessary
    %\renewcommand{\thefigure}{2}
    \centering

    \includegraphics[width=\textwidth]{#5}

    \vspace{.4cm}

    \caption{\textbf{#2.} #4}

    \label{#3} % \label works only AFTER \caption within figure environment

  \end{figure}
}
%% END MACROS SECTION
